We shall also consider the irreducible representations of the axial-symmetry group $C_{\infty v}$. This question was essentially answered already when we classified the electronic terms of a diatomic molecule, whose symmetry is precisely $C_{\infty v}$ when its two atoms are different. Two one-dimensional representations correspond to the terms $0^+$ and $0^-$ (the terms with $\Omega=0$): the identity representation $A_1$, and the representation $A_2$, in which the basis function is invariant under every rotation but changes sign under reflections in the planes $\sigma_v$. The doubly degenerate terms with $\Omega=1,2,\ldots$ correspond instead to two-dimensional representations, denoted by $E_1,E_2,\ldots$ Under a rotation through $\varphi$ about the axis their basis functions are multiplied by $e^{\pm i\Omega\varphi}$, while reflection in a plane $\sigma_v$ interchanges the two functions. The characters of all these representations are
\begin{equation}
\begin{array}{c|ccc}
 C_{\infty v}&E&2C(\varphi)&\infty\sigma_v\\ \hline
 A_1&1&1&1\\
 A_2&1&1&-1\\
 E_k&2&2\cos k\varphi&0
\end{array}
\tag{98.5}
\end{equation}
The irreducible representations of $D_{\infty h}=C_{\infty v}\times C_i$ follow directly from those of $C_{\infty v}$; they correspond to the classification of the terms of a diatomic molecule with identical nuclei.

